\section{Цель работы}
Приобретение практических навыков в:
\begin{itemize}
\item освоении принципов работы с файловыми системами;
\item обеспечении обмена данными между процессами посредством технологии <<File mapping>>.
\end{itemize}

\section{Задание}
Составить и отладить программу, осуществляющую работу с процессами и взаимодействие между ними в одной из двух операционных систем. Основной процесс должен создать для решения задачи один или несколько дочерних процессов. Взаимодействие между процессами осуществляется через системные сигналы или события и (или) через отображаемые файлы (memory-mapped files). Необходимо обрабатывать системные ошибки, которые могут возникнуть в результате работы~\cite{task}.

Вариант выбирается так же, как в лабораторной работе №2, то есть это вариант 5. Первая строка, введённая пользователем в консоль родительского процесса, содержит имя выходного файла. Родитель передаёт его дочернему процессу при запуске. Последующие строки содержат по одному числу. Дочерний процесс проверяет число на простоту: составные числа записываются в выходной файл, а при получении отрицательного или простого числа оба процесса завершают работу. Отличие от лабораторной работы №2 состоит в способе обмена: вместо неименованных каналов команды и ответы передаются через отображаемый файл.

\subsection{Принятые уточнения}
Реализация выполнена на C++17 на основе того же платформенного интерфейса \texttt{os}, что и лабораторная работа №2. Создание процессов и запись выходного файла выполняются через интерфейс \texttt{os}, а отображение файла в память и уведомления между процессами --- средствами операционной системы в отдельном модуле \texttt{ipc}. Реализация поддерживает Windows и Linux. Общая логика разбора числа и проверки простоты используется из лабораторной работы №2 без копирования.

Условия, не определённые заданием, выбраны такими же, как в лабораторной работе №2. Команда содержит одно целое число в диапазоне типа \texttt{int}. Числа 0 и 1 не являются ни простыми, ни составными: они пропускаются без завершения обработки. Некорректная строка отклоняется, после чего допускается ввод следующей команды. Конец стандартного ввода завершает обмен. Выходной файл создаётся либо обнуляется при запуске дочернего процесса.

Дополнительное ограничение связано с общим буфером: строка команды должна быть короче 256 байт. Более длинная строка считается некорректным вводом. В лабораторной работе №2 такого ограничения не было, так как строка передавалась через канал.
